For each experiment, a minimum of two sets of KOSs was required:
One set serving as a reference and baseline, with the Arduino receiving unadapted targets, and the second set used for testing target adaptation.
All experiments used a \SI{60}{\second} buffer and \SI{15}{\second} intervals between updates, with the first update at \SI{60}{\second}.
The adapter trained on prediction windows of 3, 5, 10, 15, and 25 steps.
The true target, with values depending on the experiment, was updated every \SI{5}{\second}, following a predefined random sequence.
This sequence amounts to \SI{5}{\minute} of deployment time and is repeated according to the full experiment duration.
The signals corresponding to each \SI{5}{\minute} sequence are recorded and saved for evaluation.
This way, the evolution of the performance metrics can be studied, and a fair comparison between the reference controller and adapted controller is ensured.
\Cref{fig:test-setup} shows the general experimental setup.

\begin{figure}[htb]
    \centering
    \missingfigure[figwidth=0.9\linewidth, figheight=0.6\linewidth]{A wide shot of the laptop connected to arduino and robot, as during experiments.  Annotations for the robot power supply, and the serial connection (+ power) for the Arduino}
    \caption{The test setup.}
    \label{fig:test-setup}
\end{figure}

\subsection{Sign Conventions and Unit Clarification}\label{subsec:sign-unit}
Before the experiments and results can be properly understood, it is necessary to clarify certain sign and unit conventions.
Most importantly, the contactless magnetic sensor denotes a clockwise rotation (as viewed from the top) w.r.t.~the datum as a positive angle, and likewise, a counterclockwise rotation as a decrement in the angle.\todo{cite arduino library}  % cite arduino library
In practice, this means that a positive control input from the servo (i.e.~a positive angle) results in a decrement in $\theta$.
Because the actuation arm's calibration point is set when the servo is at its minimum output, the operational range of the actuation arm is entirely in the negative domain.
\Cref{fig:schematic} illustrates the servo arm and actuation arm in a configuration representative of a typical scenario during the experiments, with the signs of the angles indicated.

\begin{figure}[htb]
    \centering
    \def\svgwidth{0.5\linewidth}
    \footnotesize\selectfont\input{Thesis/figures/robot/schematic.pdf_tex}
    \caption{Nominal signs of the angles for the servo and actuation arm during operation.}
    \label{fig:schematic}
\end{figure}

In addition, The \emph{Servo} Arduino library uses the \texttt{writeMicroseconds(\emph{value})} function to write to the servo, which maps a servo's range to a value between 1000 and \SI{2000}{\micro\second} depending on the servo's manufacturer.
As such, \texttt{writeMicroseconds} provides more `resolution' than the common \texttt{write(\emph{angle})}, which works with integer values of angles in degrees. \todo{cite Servo library} % TODO: cite
Note that the term \emph{microseconds} in this context does not refer in any way to the angular unit of measurement \emph{arcseconds}.
Instead, it relates to the pulse width in microseconds of the signal that the Arduino sends to the servo motor.
In this case, the servo's minimum position corresponds to a value of \SI{540}{\micro\second} and the maximum mechanical limit (as discussed in \Cref{subsec:mechanical-details}) corresponds to approximately \SI{1800}{\micro\second}.
It might thus be confusing to see the control input, the angle of the servo arm, denoted in a unit of time.
For convenience, the \SI{540}{\micro\second} value relates to 0 degrees, and the \SI{1800}{\micro\second} limit corresponds to 125 degrees.

Finally, during the experiments, the Arduino initialized a \emph{counter} value to include in the tuple sent to the computer.
This counter starts at 0 and simply increments by 1 with each iteration of the Arduino's loop.
Since the interactions between the Arduino and computer, as well as between the several threads on the computer, are independent of timing (i.e., the modules do not have to ``wait for each other''), the counter's value is essential to accurately plot the several signals recorded during the experiments.
As a result, the counter itself does not measure elapsed time directly, making it unitless.
However, it can be used as a proxy for time, since all modules aim to operate at \SI{50}{\hertz}.
This means that an increment of 1 in the counter corresponds to an estimated \SI{0.02}{\second}, allowing for the approximation of elapsed time based on the counter value.

\subsection{Performance Metrics}\label{subsec:metrics}
The performance metrics used to evaluate the \SI{5}{\minute} sequences are the mean rise time, mean settling time, mean overshoot, mean absolute error, integrated absolute error, mean absolute velocity, normalized total variance, and mean control action.

The mean rise time is calculated as the mean of all the rise times corresponding to each target in the \SI{5}{\minute} sequence.
The rise time is defined as the time it takes the robot to come within 10\% of the target state, as calculated from the initial position.

Similarly, the settling time is the average of the settling times for each target.
The robot state is considered settled when it remains within 5\% of the target state, as calculated from the initial position.

The mean overshoot is the average of the maximum overshoots observed for each target.

The mean absolute error represents the average difference between the robot’s angle $\theta$ and the true target $\theta_d$, while the integrated absolute error captures the cumulative difference over time.

The mean absolute velocity refers to the measurements of $\dot{\theta}$.
This metric is meant to provide a measure of the ``chaoticness'' of the robot's response signal.

The normalized total variance is calculated as the sum of the absolute of all increments and decrements in the control action, normalized over time\footnote{Actually, it is normalized over the \emph{count}, see \Cref{subsec:sign-unit}.}.
This, combined with the mean control action, aims to quantify the controller's ``effort''.